% =====================================================================
% tesi/capitoli/24-strumenti.tex
% =====================================================================
% copre: docs/22-strumenti.md
% copre: pokemon-gen12-gen3-bridge-original-hardware/DATA-FORMATS_Gen1-Gen2-Gen3.md#7-codifica-dei-caratteri
% ---------------------------------------------------------------------

\chapter{Gli strumenti, e il principio che li governa}
\label{cap:strumenti}

Gli strumenti prodotti da questo lavoro sono pochi e deliberatamente piccoli, e rispondono a un principio che vale enunciare prima di descriverli: in una catena di lavorazione che mescola codice e giudizio, si spinge sul codice deterministico tutto ciò che non richiede comprensione semantica, e se ne conserva l'esito come stato ispezionabile su disco anziché come risposta in una conversazione.

I benefici di questo principio sono tre e sono verificabili. La riproducibilità, perché rieseguendo il programma si ottiene il medesimo risultato. L'economia, perché il lavoro deterministico costa tempo di calcolo locale e non risorse di ragionamento. E l'ispezionabilità, perché gli stati intermedi sono file leggibili che si possono correggere a mano senza rieseguire l'intera catena.

\section{Il presupposto comune: i disassemblati in locale}

Due dei tre strumenti principali leggono dai disassemblati, i quali non risiedono nel repository e non devono risiedervi. Si ottengono con cloni superficiali, che sono sufficienti perché serve soltanto lo stato corrente.

\begin{verbatim}
git clone --depth 1 https://github.com/pret/pokecrystal
git clone --depth 1 https://github.com/pret/pokeemerald
git clone --depth 1 https://github.com/pret/pokered
\end{verbatim}

Conviene conservarli fuori dal repository. Il vantaggio del clone rispetto alla lettura attraverso il web è che consente di cercare nel codice anziché leggere descrizioni, ed è precisamente in questo modo che sono state individuate le quattro correzioni documentate nella premessa. Vale generalizzare l'osservazione, perché è uno dei risultati metodologici di questo lavoro: su dati di gioco il clone superficiale è il primo passo di qualunque verifica, non l'ultimo.

\section{Il generatore delle tabelle di codifica}

Il primo strumento genera le tabelle di codifica dei caratteri leggendo i file di definizione dei disassemblati, e produce tre file. Il primo è la tabella delle generazioni 1 e 2, il secondo quella della generazione 3, ciascuna con i caratteri stampabili separati dai token di controllo e con l'indicazione del commit di provenienza. Il terzo è la traduzione diretta fra i due spazi di codifica, che è il file che un convertitore impiega effettivamente, e che elenca separatamente i caratteri privi di destinazione.

Le cifre della corsa corrente sono duecento caratteri stampabili per le generazioni 1 e 2, duecentoquarantotto per la generazione 3, centoquarantasette corrispondenze e cinquantatré caratteri privi di destinazione. Quest'ultimo numero non è un residuo ma un dato di progetto, per la ragione esposta nel capitolo~\ref{cap:testo}: che cosa fare di un nome che contenga quei caratteri è una decisione, e il file la rende visibile.

La proprietà rilevante di questo strumento non è ciò che produce ma ciò che rifiuta di produrre. Prima di scrivere verifica un insieme di valori di controllo, cioè i byte del terminatore, dello spazio, delle lettere estreme dei due alfabeti e delle cifre estreme in entrambe le tabelle, e qualora anche uno solo non corrisponda non scrive nulla e riferisce lo scostamento. Se il formato dei file a monte cambiasse, il fallimento sarebbe rumoroso anziché silenzioso, che è esattamente il contrario del comportamento della trascrizione manuale il cui esito è documentato nel capitolo~\ref{cap:testo}.

\section{Lo strumento di diagnosi dello zaino}

Il secondo strumento legge un salvataggio di generazione 3 per diagnosticare lo zaino, e non scrive nulla: legge, valida e riferisce. La sequenza è quella descritta nei capitoli~\ref{cap:integrita} e~\ref{cap:cifratura}, cioè individua i due slot, valida la firma e il checksum di ciascuna sezione, seleziona lo slot più recente secondo la regola del contatore, ricompone il blocco di salvataggio dalle sezioni, identifica il gioco, legge la chiave di cifratura al proprio offset e smaschera le quantità dello zaino lasciando in chiaro quelle del deposito.

L'identificazione del gioco non si fida di un parametro: lo verifica. La ragione proviene da un caso reale documentato in una discussione della comunità, in cui un editor aveva identificato un salvataggio di Smeraldo come Rubino o Zaffiro e gli oggetti sono finiti negli slot sbagliati, il che è esattamente l'effetto dell'applicazione della maschera errata \cite{projectpokemon}. Lo strumento assegna un punteggio a ciascun candidato sulla base di prove indipendenti, cioè la plausibilità del conteggio della squadra al proprio offset, il denaro smascherato entro il tetto, e la chiave dedotta da uno slot vuoto, e stampa le prove di tutti i candidati anziché dichiarare soltanto il vincitore. Qualora due candidati pareggino si arresta e chiede l'indicazione manuale, perché un'identificazione ambigua è un'informazione e non un dettaglio da risolvere per estrazione.

Riferisce poi le anomalie che sa riconoscere: un identificatore di oggetto fuori intervallo, una quantità nulla su uno slot occupato, una quantità oltre il tetto della tasca, uno slot popolato dopo un buco, e un identificatore duplicato nella medesima tasca. Sono cinque classi che coprono i sintomi tipici di uno zaino alterato da codici di alterazione.

A queste si aggiunge un controllo di categoria, introdotto per la casistica concreta esposta nel capitolo~\ref{cap:smeraldo}, cioè quella in cui oggetti ordinari risultano sostituiti da Poké Ball a partire da una certa posizione. Il controllo verifica che l'identificatore di ciascun oggetto appartenga all'intervallo di categoria coerente con la tasca in cui risiede, impiegando gli intervalli verificati sul disassemblato, e riferisce se la difformità possieda un punto di inizio, cioè se esista una posizione a partire dalla quale ogni slot successivo sia difforme. La presenza di un punto di inizio è essa stessa un dato diagnostico, perché distingue una corruzione progressiva da un'alterazione puntuale.

Lo strumento possiede due verifiche incrociate che valgono più del resto, e la loro natura è quella descritta nel capitolo~\ref{cap:collaudo}: due vie indipendenti per il medesimo valore. La prima è che il denaro, smascherato con la medesima chiave, deve stare sotto il proprio tetto: in caso contrario la chiave o l'offset sono errati, e lo strumento lo dichiara anziché produrre numeri privi di senso. La seconda è che uno slot vuoto in una tasca mascherata contiene la chiave stessa, dunque la chiave si ricava per una seconda via e le due si confrontano.

Sul lato dei titoli della coppia Rosso Fuoco e Verde Foglia lo strumento è esplicito sui propri limiti: legge la chiave al proprio offset e la riferisce, ma non pretende di elencare le tasche, perché i loro offset non sono stati verificati sul disassemblato di quel gioco in questa revisione. La dichiarazione di un limite è parte del contratto dello strumento, non un'omissione.

\section{Il pacchetto del ponte}

Il pacchetto non è uno strumento ma la prima parte del software vero, e copre gli strati dal primo al terzo della stratificazione descritta nel capitolo~\ref{cap:ponte}, cioè i dati generati, i modelli e i lettori e scrittori. Non possiede dipendenze esterne.

I primitivi risiedono nel modulo che ospita la conoscenza non appartenente a una singola generazione: interi a sedici e ventiquattro bit big-endian, scomposizione dei nibble dei valori individuali con la derivazione del quinto, byte dei punti potenza diviso secondo lo schema a sei più due bit, e il pattern di valori che nella generazione 2 significa lucentezza. I lettori e scrittori risiedono in un modulo per generazione, perché fra le prime due esiste un riordino e non un'estensione, come il capitolo~\ref{cap:gen2} stabilisce. La transcodifica legge le tabelle generate e non contiene alcun valore scritto a mano.

Vale registrare una circostanza sulla prima esecuzione delle prove, perché ripete la lezione dello strumento precedente applicandola alle prove stesse: una prova è fallita, e il difetto era nella prova e non nel codice, perché il byte di una lettera era stato scritto a mano con il valore sbagliato. La prova ricava adesso i byte dalla tabella anziché dichiararli, che è la medesima disciplina del generare anziché trascrivere applicata al collaudo.

\section{Gli strumenti di infrastruttura}

Esistono infine gli strumenti che servono al repository e non ai giochi, e la loro presenza in questo capitolo è giustificata dal fatto che ciascuno attua meccanicamente una convenzione che altrimenti resterebbe dichiarata e non verificata.

Il primo attua la convenzione di formattazione dei documenti, cioè un paragrafo per riga sorgente, e dispone di una modalità di sola verifica per il controllo che precede una registrazione delle modifiche. Il secondo verifica che i comandi di shell dentro i blocchi di codice non siano spezzati su più righe, perché il primo per contratto non tocca il contenuto dei blocchi di codice e dunque un comando spezzato in quella sede non viene corretto da nessuno.

A questi si aggiungono tre strumenti tipografici che attuano le convenzioni sull'ortografia: il primo converte le vocali accentate scritte con l'apostrofo nella loro forma corretta, distinguendo l'accento acuto dal grave secondo la grammatica; il secondo interviene sugli accenti mancanti del tutto, e non pretende di chiudere il problema da solo, perché converte le sole forme che privi di accento non sarebbero parole e riporta come indecidibili quelle che richiedono di comprendere la frase; il terzo normalizza i trattini, eliminando quelli lunghi in ogni loro variante, incluso il segno meno matematico, che in un testo tecnico è il più insidioso perché è un operatore e chi copia una formula che lo contiene ottiene un carattere che nessun compilatore accetta.

Una convenzione dichiarata e non verificata è esattamente ciò che ha prodotto le tre grafie diverse che questi strumenti hanno dovuto sanare, e questo è l'argomento che ne giustifica l'esistenza.

\section{Il cancello prima dell'hardware}

Tre dei cinque casi di studio convergono, presto o tardi, sulla medesima operazione: prendere un file di salvataggio e collocarlo su un supporto fisico, cioè una cartuccia attraverso il lettore oppure la scheda di memoria della console. È l'unica operazione irreversibile del progetto, e per questa ragione esiste uno strumento unico anziché tre procedure separate.

La proprietà da comprendere di questo strumento è che non scrive, e non scriverà finché l'hardware non sia presente e collaudato. Ciò che compie è la parte che si può scrivere oggi e che evita i danni. Esamina il file: la dimensione confrontata con quelle che possiedono un significato, l'impronta che identifica il contenuto, il numero di settori che portano la firma, il numero di settori il cui checksum torna, l'identificazione del gioco con il margine sul secondo candidato, e i due casi degeneri del file interamente a zero e del file interamente a \hx{FF}, che è lo stato di una memoria flash cancellata. Poi verifica le precondizioni della destinazione e produce il piano dei passi.

La validazione non è duplicata: lo strumento importa le funzioni di validazione dei settori, di ricostruzione dello slot e di identificazione del gioco dallo strumento di diagnosi. Un secondo esemplare della medesima formula di checksum sarebbe un secondo luogo in cui sbagliarla.

La parte più importante è il cancello dei backup, che attua in codice il vincolo normativo esposto nel capitolo~\ref{cap:perimetro}. Lo strumento pretende due backup dichiarati, verifica che i percorsi siano distinti, che i file esistano e siano leggibili, che la dimensione coincida con quella dell'originale e che le due impronte coincidano fra loro, e rifiuta anche il caso in cui i due backup risiedano sul medesimo volume, perché un guasto del supporto li perderebbe insieme. Quando una di queste condizioni manca, il piano non viene stampato affatto: è una scelta deliberata, perché una lista di passi mostrata sotto un elenco di problemi invita a eseguirla comunque.

Il piano, quando viene emesso, presenta cinque passi nell'ordine che la norma impone: dump del contenuto attuale del supporto, confronto con i backup dichiarati, scrittura, rilettura confrontata byte per byte, e infine avvio del gioco, che è un controllo diverso dal precedente e non lo sostituisce. Poi lo strumento dichiara che la scrittura non viene tentata e specifica che cosa manchi a quella destinazione: per la cartuccia il lettore e un collaudo su un esemplare sacrificabile, per la scheda di memoria la decisione su quale percorso sia quello corretto per il titolo installato.

Il collaudo è avvenuto su un salvataggio sintetico generato per l'occasione, con quattordici settori firmati e coerenti, e ha verificato i tre casi che contano: il file valido identificato correttamente con il proprio margine, il rifiuto del file interamente a zero e della dimensione non riconosciuta, e il rifiuto del piano nei tre modi in cui i backup possono risultare insufficienti, cioè uno solo, inesistente, oppure due sul medesimo volume.

\section{Il generatore del catalogo degli eventi}

Il quinto strumento genera l'inventario delle distribuzioni di evento della terza generazione discusso nel capitolo~\ref{cap:distribuzioni}, a partire dalla tabella di un verificatore di legittimità di terze parti e dall'enumerazione dei metodi di generazione contenuta nel medesimo repository. Il percorso del clone si passa sulla riga di comando, con la disciplina già adottata per il generatore delle tabelle di codifica: quel clone non è una dipendenza del presente lavoro e non viene acquisito dallo strumento.

La ragione per cui il catalogo viene generato anziché trascritto coincide con quella delle tabelle di codifica, ed è la ragione più generale per cui questo lavoro preferisce generare. I dati risiedono nel codice di chi verifica e non in un documento, poiché, come il commento della tabella stessa dichiara, i dati di quella generazione non sono mai stati conservati in forma binaria uniforme e sono scritti a mano voce per voce; trascriverli garantirebbe la divergenza delle due copie alla prima correzione a monte. Lo strumento dispone di una modalità di sola verifica che stabilisce se il catalogo sia ancora allineato alla fonte, ed è il controllo da eseguire quando il clone viene aggiornato.

Ciò che lo strumento non compie va dichiarato con la medesima cura: non giudica la legittimità di alcun esemplare e non ricostruisce alcun metodo, riporta quanto la fonte dichiara. Le date degli eventi compaiono soltanto dove la fonte le riporta, e dove non le riporta restano assenti anziché essere congetturate.

\section{Il lettore di canali di comunità}

Il sesto strumento legge la cronologia di un canale di conversazione attraverso un account dedicato all'automazione, secondo la terza delle tre vie discusse nel capitolo~\ref{cap:perimetro}, con impaginazione, un cursore che consente di leggere il solo incremento fra due esecuzioni, e i medesimi filtri dello strumento che digerisce le trascrizioni manuali, affinché il lavoro non conosca due grammatiche di filtro.

Le difese contro le condizioni reali di un canale affollato sono state aggiunte in un secondo momento, e la circostanza merita di essere esposta perché costituisce un caso di studio in miniatura del limite di ogni collaudo. La prima esecuzione riuscita era avvenuta su un apparato di prova contenente cinque messaggi, e aveva dato esito positivo su ogni passo; ma un canale con cinque messaggi non esercita l'impaginazione, non raggiunge alcun limite di frequenza, non produce guasti di rete, non contiene discussioni annidate né contenuti non testuali, e non ospita testo di terzi capace di interferire con la struttura del documento prodotto. Un collaudo riuscito, in altri termini, misura anche ciò che l'apparato di prova era in grado di sollecitare, e quel secondo dato non compare nel suo esito.

Le difese introdotte sono sette e ciascuna corrisponde a una di quelle assenze. Il limite di frequenza è governato in modo reattivo, attendendo quanto il servizio dichiara nel proprio rifiuto, e in modo preventivo, leggendo a ogni risposta le intestazioni con cui il servizio annuncia le richieste residue nella finestra corrente: la seconda modalità evita il rifiuto anziché reagirvi. I guasti transitori, ossia gli errori di rete e le risposte di errore del servizio, non interrompono una lettura a metà ma provocano una ripetizione con attesa raddoppiata fino a un tetto, e l'abbandono definitivo riferisce l'ultimo esito osservato. Le discussioni e i post di forum, che nell'interfaccia appaiono annidati e nell'API sono canali con identificativo proprio, vengono elencati interrogando l'endpoint dedicato, poiché la loro omissione farebbe apparire vuoto un canale ricco. Il contenuto non testuale, ossia i blocchi incorporati e la citazione del messaggio cui una risposta si riferisce, entra nella resa e nei filtri: ignorarlo renderebbe vuoti messaggi che portano contenuto, e li farebbe scartare proprio dal filtro di lunghezza. Il testo di terzi che apre con un carattere di intestazione viene protetto, affinché non forgi la struttura del documento, con l'eccezione dei blocchi di codice recintati dove quel carattere appartiene al linguaggio. La data impiegata come filtro viene verificata nella forma, poiché una data scritta in una notazione diversa passerebbe senza errore e scarterebbe la totalità o nulla. E la scrittura può aggiungersi in coda a un documento esistente anziché sovrascriverlo.

La lezione generale vale enunciarla, poiché si applica a ogni collaudo di questo lavoro: l'esito di una prova va letto insieme all'inventario di ciò che l'apparato di prova non poteva sollecitare, e quel secondo elenco va scritto, perché non emerge da sé.

La parte che merita esposizione è il presidio, poiché esemplifica un principio applicabile oltre il caso. Lo strumento trasmette sempre l'intestazione di autorizzazione nella forma destinata agli account di automazione e, prima di qualunque lettura, verifica che l'account autenticato sia dichiarato tale, arrestandosi con la propria ragione in caso contrario: un token personale inserito per errore nella variabile d'ambiente non produce una lettura riuscita ma un rifiuto. Il principio è che una distinzione normativa, per essere effettiva, deve essere resa meccanica nel punto in cui potrebbe essere violata per distrazione: dichiararla nella documentazione la rende conoscibile, verificarla nel codice la rende operante.

La prova del presidio è un controllo negativo, cioè una verifica che fallisce se il presidio viene rimosso, e appartiene a una suite che esercita anche impaginazione, cursore, attesa dopo un rifiuto per eccesso di frequenza, controllo locale degli identificativi e filtri contro un trasporto simulato.

Il flusso verso il servizio è stato in seguito eseguito su un apparato di prova di proprietà di chi conduce il lavoro, e ha funzionato in ciascuno dei propri passi. Ciò che resta non provato sul servizio va dichiarato con la medesima precisione, ed è quanto un apparato di prova non può esercitare: l'impaginazione su una cronologia più estesa del massimo restituito da una singola richiesta, e l'attesa successiva a un rifiuto per eccesso di frequenza. Entrambe sono provate contro il trasporto simulato, nessuna delle due è stata osservata sul campo, e la distinzione fra i due stati è quella che il capitolo~\ref{cap:collaudo} impone all'intero lavoro.

Dalla medesima esecuzione provengono due osservazioni che appartengono al genere di cui il capitolo~\ref{cap:collaudo} raccomanda la registrazione, poiché nessuna delle due è deducibile dalla documentazione. La prima è che fra i messaggi restituiti compare anche la riga con cui il servizio annuncia l'ingresso dell'account automatico nel canale, essendo per il servizio un messaggio come gli altri. La seconda è che il messaggio di errore restituito quando l'identificativo non è un numero non nomina il campo difettoso: è la ragione per cui la verifica dell'identificativo è stata spostata nel programma, ed è un caso particolare del principio esposto poco sopra, cioè che un presidio va collocato nel punto in cui l'errore si produce e non in quello in cui si manifesta.

Una scelta di progetto va infine motivata, poiché il materiale che ha aperto la questione raccomandava una soluzione diversa, cioè l'interposizione di un server conforme al protocollo per l'integrazione di strumenti in un agente. In un agente residente che debba invocare quello strumento nel corso di una conversazione quella è la scelta corretta; qui il compito è deterministico, cioè leggere una fonte e trasferirne la sintesi nel registro, e il criterio esposto nel capitolo~\ref{cap:collaudo} prescrive di collocare il lavoro deterministico nel codice anziché nel modello. Un programma fondato sulla sola libreria standard assolve quel compito senza introdurre una dipendenza da un ambiente di esecuzione ulteriore, un pacchetto di terze parti a cui affidare una credenziale, e uno strato di protocollo interposto fra il chiamante e una richiesta HTTP.

\section{Uno strumento di terze parti, e la divisione del lavoro con quello proprio}

Accanto agli strumenti scritti per questo lavoro se ne impiega uno di terze parti, maturo e diffuso, che esporta la cronologia di un canale di conversazione in più formati e scarica gli allegati. La sua presenza in questo capitolo non è un elenco di dipendenze ma un caso di studio sulla divisione del lavoro fra uno strumento generale e uno specifico, e su come si documenta una procedura che non è codice.

La divisione è la seguente, e discende dalle rispettive proprietà anziché da una preferenza. Lo strumento di terze parti è superiore su tre fronti: porta in locale gli allegati, che in un canale tecnico sono schemi e schermate e non ornamento; produce una resa che replica l'interfaccia originale, ed è quindi il formato in cui una discussione lunga si rilegge; e dispone di un'interfaccia grafica, comoda quando occorre percorrere molti canali per stabilire quali contengano qualcosa. Lo strumento proprio è superiore su tre altri: possiede il cursore che legge il solo incremento fra due esecuzioni senza che alcun identificativo vada gestito a mano, consegna il risultato già filtrato nella convenzione del presente lavoro, e reca il presidio sul tipo di credenziale. Ne segue che il primo compie le esportazioni ampie e infrequenti e il secondo gli aggiornamenti incrementali, e che il ponte fra i due preesisteva, poiché il convertitore descritto sopra era stato scritto per digerire il formato del primo.

Sulla credenziale con cui quello strumento opera il presente lavoro ha compiuto una scelta che appartiene al capitolo~\ref{cap:perimetro} e che va nominata anche qui, poiché determina la portata di ciò che si può leggere: dove la via che passa da un account dedicato all'automazione non è disponibile, per assenza del consenso di chi amministra il canale, si impiega la credenziale personale, con la consapevolezza documentata delle sue implicazioni e con una cadenza di poche esecuzioni all'anno. La decisione è registrata con i suoi termini reali fra le decisioni del progetto, ed è esposta là anziché attenuata qui.

Un'ultima osservazione riguarda la forma della documentazione e vale oltre il caso. La procedura d'uso di uno strumento di terze parti non è codice e non ha alcuna verifica automatica che ne rilevi l'obsolescenza: né una suite di prove, né un controllo di copertura, né un compilatore. Ne segue che va scritta nel repository con la medesima cura di un capitolo, con i percorsi, i nomi esatti dei file e la ragione di ciascuna opzione, poiché una procedura che risieda in una conversazione è perduta alla sessione successiva, e una che ometta il perché di un'opzione verrà eseguita senza quell'opzione dal primo che la riprenda.

\section{Lo strumento che confronta il proprio sapere con quello di un altro}

Gli strumenti descritti fin qui producono un dato o ne verificano la coerenza interna. L'ultimo appartiene a una categoria diversa e vale descriverlo per essa, poiché il criterio che lo giustifica è riusabile: mette a confronto ciò che questo lavoro ha stabilito con ciò che un'implementazione indipendente della stessa materia dichiara, e lo fa in modo ripetibile anziché aneddotico.

L'occasione era una domanda di metodo del capitolo~\ref{cap:provenienza}, cioè se le due vie che compongono un esemplare da evento producano lo stesso dato. La forma della risposta poteva essere una lettura comparata, che è ciò che si fa di solito, oppure un programma. La scelta del programma poggia su una ragione che vale enunciare: una lettura comparata produce un giudizio non riproducibile e invecchia in silenzio, mentre un confronto meccanico si può rieseguire quando una delle due parti cambia, e il suo esito è un numero anziché un'impressione.

Lo strumento esegue cinque confronti e li ordina per durezza. I primi quattro mettono a paragone due dichiarazioni scritte in codice: una tabella di permutazione, una tabella di codifica dei caratteri sulla loro intersezione, un vocabolario di metodi come confronto fra insiemi, e un inventario congiunto sulla sola chiave presente in entrambe le fonti. Il quinto è di natura diversa e vale più degli altri, perché non paragona due descrizioni ma esegue il codice di questo lavoro sui dati conservati dell'altra parte e verifica che ne riproduca gli esiti.

Tre lezioni di implementazione meritano di essere registrate, poiché nessuna era prevedibile e ciascuna ha un valore generale.

La prima riguarda il recupero del materiale. Il filtro che riconosce un riferimento di importazione rifiuta ciò che contiene caratteri non ammessi in un percorso, e questo non è eccesso di prudenza: senza di esso l'espressione regolare aggancia anche le stringhe contenute nei commenti, e la prima versione del programma ha inviato al servizio remoto un percorso che conteneva un intero blocco di codice. Il principio è che un'espressione regolare applicata a un linguaggio non è un analizzatore sintattico, e va corredata di un controllo di plausibilità sul risultato.

La seconda riguarda la codifica dell'uscita. I nomi degli allenatori giapponesi non passano dalla codifica predefinita del terminale del sistema impiegato, e in assenza di una riconfigurazione esplicita il programma termina con errore alla stampa finale, dopo avere eseguito correttamente ogni confronto. È il modo peggiore di fallire, perché il lavoro è fatto e il risultato è perduto, e la regola che ne discende è che la presentazione va resa robusta indipendentemente dal calcolo.

La terza riguarda l'interpretazione dell'esito, ed è la più importante. Il confronto ha inizialmente segnalato sedici assenze nell'inventario della controparte, che sembravano lacune sue e non lo erano: quella implementazione conserva il proprio inventario in due luoghi, un albero di moduli e un corpus caricato a tempo di esecuzione, e la nostra lettura ne consultava uno solo. Il criterio generale è che un'assenza rilevata da un confronto è un'ipotesi sulla controparte e va verificata sulla propria procedura di lettura prima di essere attribuita a essa.

L'esito, riportato nel capitolo~\ref{cap:distribuzioni}, è che le due parti concordano sui dati e che due difetti sono stati individuati nella controparte, dei quali uno con conseguenza operativa. Vale osservare che il secondo esito non era l'obiettivo dello strumento: un confronto costruito per rispondere a una domanda su di sé ha prodotto, come effetto collaterale, informazione sull'altro. È la proprietà per cui questo genere di strumento merita di essere costruito anziché sostituito da una lettura.

\section{Che cosa manca}

Lo strumento successivo in ordine naturale è il generatore della tabella dagli indici interni della generazione 1 ai numeri nazionali, che oggi non esiste e che va costruito con il medesimo metodo, cioè leggendolo dal disassemblato e verificandolo con valori di controllo. È l'ultimo dato costante del progetto che sarebbe tentante trascrivere a mano, ed è anche quello in cui un errore silenzioso produrrebbe il danno maggiore, perché scambierebbe una specie con un'altra.
